% conclusion.tex

\section{CONCLUSIONS} \label{sec:conclusions}

The laboratory confirms that SSH honeypot sessions can be treated effectively
as contextual token sequences rather than as isolated command names.
Tokenization is a central design choice: before preprocessing, long encoded
strings create pathological sequences, while truncating words above 30
characters makes both BERT and UnixCoder usable under the 512-token limit.

In the NER stage, pre-training is clearly beneficial. Naked BERT remains well
below pre-trained BERT, and the code-oriented UnixCoder model obtains the best
overall performance, with the highest test macro F1 and session fidelity. The
frozen-layer experiments show a practical trade-off: training only part of the
encoder reduces computational cost, but full fine-tuning still provides the
best balanced performance.

Finally, inference on CyberLab logs shows how the trained model can support
security analysis. Command-level tactic profiles reveal that commands such as
\texttt{echo} and \texttt{rm} are strongly context-dependent, while fingerprint
analysis exposes persistent behaviours and high-volume December bursts that
are compatible with automated attack campaigns.
